% page 303 du fac-simile (image p-302 du scan)
% bloc 167.6 x 216.0 mm ; marge gauche 34.4 mm
\pgc{12.700mm}{0.000mm}{
\l{\cel{56}295}
\l{}
\l{\sou{konfitar}. (\sou{trans.}) Preparar (frukti) per igar li sejornar en}
\l{\cel{3}liquoro, en siropo, qua penetras en li, qua konservas li.}
\l{\cel{3}- DFIS.}
\l{}
\l{\sou{konfliktar}. (\sou{netrans., kun)}. I. 1.Esar en interlukto, inter-}
\l{\cel{3}kombato (\sou{aludante personi qui interbatas}). 2. (metaf.)}
\l{\cel{3}(\sou{pri la interesti, la pasioni}).- II. Interkontestar, (aludante}
\l{\cel{3}du\cel{1}\sou{korpi politikala, judiciala, e c., singla de qui volas havar}}
\l{\cel{3}\sou{kom suan, yuro, resortiso, fako, e}\cel{1}c.) - DEFIRS.}
\l{}
\l{\sou{konfluanta}. Di qua la elementi interjuntas. - (medic.) Pustuleti}
\l{\cel{3}\sou{konfluanta}\cel{1}: tante interproxima ke li intertushas ed inter-}
\l{\cel{3}mixesas. - (\sou{bot.})\cel{1}\sou{Folii konfluanta}\cel{1}: di qui la stipiti interjun}
\l{\cel{3}tas ed intermixas.}
\l{}
\l{\sou{konforma}. (\sou{ad)}. I. Di qua la formo koincidas precize kun olta}
\l{\cel{3}di altra objekto quan onu egardas kom tipo. - II. Qua koincidas}
\l{\cel{3}precize kun to quo imperesas da ulu, o kun to quo igesas}
\l{\cel{3}necesa da ulo. - DEFIS.}
\l{}
\l{konfrontar. (\sou{trans.,}\cel{1}ad).I. Situar facio vers facio (du o}
\l{\cel{3}plura personi da qui la afirmi esas interkontrea) por dicer-}
\l{\cel{3}nar lo exakta. - II. Inter-apud-pozar (texti, dokumenti)}
\l{\cel{3}por interkomparar li. - DEFIS.}
\l{}
\l{konfundar. (\sou{trans., ad, kun}). I.\cel{1}\sou{Konfundar kun}\cel{1}: Facar pelmelo}
\l{\cel{3}pro distingo mentala qua esas mala o nesuficanta, e mixas ici}
\l{\cel{3}kun iti (kom ex. \textquotedbl{}Il konfundas omno en sua odio, la kulpozi}
\l{\cel{3}kun l\textquotesingle{}inocenti, la benigni kun la maligni\textquotedbl{}) - II.\cel{1}\sou{Konfundar}\cel{1}ad}
\l{\cel{3}Transportar (mentale) erore la qualesi o la individueso di ulu}
\l{\cel{3}od ulo, ad altra (kom ex.: \textquotedbl{}Vu konfundis manekino ad homo. Ne}
\l{\cel{3}konfundez lampiro ad lanterne\textquotedbl{}) DEFIRS.}
\l{konfuza. Di qua la elementi esas intermixita tale ke li cesis}
\l{\cel{3}esar dicernebla e distingebla. - DEFIRS.}
\l{konglomerato. (mineral.) Bloko de fragmenti minerala aglomerita.}
\l{\cel{3}- DEFIRS.}
\l{kongro. (\sou{zool}.) Marfisho ek la sama familio kam la mureno.}
\l{\cel{3}- L.\cel{1}\sou{muraena conger}. - EFIS.}
\l{kongregaciono. I. (a) Kompanio de sacerdoti qui su obligis per}
\l{\cel{3}vovi, per regulo komuna, agnoskita da la Santa\cel{2}Siejo ;}
\l{\cel{3}(b) Kompanio de kleriki qui, aganta segun la regulo komuna}
\l{\cel{3}di la ordeno a qua li apartenas, su obligas per regulo spe-}
\l{\cel{3}cala. - II. Asociitaro de nekleriki, qui interunionas sub}
\l{\cel{3}la nomo di santo, por la praktikado di pieso. - III. Komitato}
\l{\cel{3}de kardinali, de kleriki, establisita da papo, en Roman, por}
\l{\cel{3}studiar ula aferi. - DEFIRS.}
\l{}
\l{kongresar. (\sou{netrans.)}\cel{1}I. (\sou{aludante ciencisti, publicisti, politik-}}
}
